\section{Preprocessing and evaluation protocol}
\label{protocol}

This section turns the decisions of Section~\ref{data} into a pipeline that every model shares.
The guiding rule comes from the course: clean first, then split, then fit, and never let a
validation or test row influence a fitted statistic.

\subsection{Cleaning}

Dropping the exact duplicates leaves 17,599 ads, 856 of them fraudulent. Each ad is then given a
\emph{description group}: the hash of its normalised description text. The 17,599 ads fall into
14,592 groups, and ads in the same group are treated as one unit when splitting.

\subsection{Grouped train, validation and test splits}

\texttt{StratifiedGroupKFold} on the description group gives a 60/20/20 split that keeps the fraud
rate in every part: 10,559 training, 3,520 validation and 3,520 test ads, with 514, 171 and 171
frauds. A check confirms that no description group appears in two parts. The three parts have
separate roles: models are fitted on the training rows, all choices (hyper-parameters, epochs,
thresholds) are made on the validation rows, and the test rows are only used for the final score.

One split gives one number with no error bar, so we repeat the whole procedure for five split
seeds (0 to 4). The row indices of each split are saved together with a fingerprint of the cleaned
data, so every notebook and every machine uses exactly the same rows: the laptop for the classical
models and the GPU server for the deep ones.

\subsection{Feature construction}

The five text fields are joined into one text column, and one \texttt{ColumnTransformer} turns
each ad into a sparse model matrix with five blocks (Table~\ref{tab:features}). The transformer is
part of each model's pipeline, so it is fitted on the training rows only.

\begin{table}[tb]
\centering
\small
\begin{tabular}{l>{\raggedright\arraybackslash}p{0.5\linewidth}r}
\toprule
\textbf{Block} & \textbf{Transformer} & \textbf{Columns} \\
\midrule
Text & TF-IDF (sublinear term frequency, word 1- and 2-grams, at least 3 documents) or BM25 & 20,000 \\
Nominal & One-hot of employment type, industry, function and country; categories under 30 ads share one column & 110 \\
Ordinal & Required experience and education as ordered levels & 2 \\
Numeric & Description and title length, exclamation marks, share of capitals; scaled & 4 \\
Binary & Telecommuting, logo, questions and 9 engineered flags; near-constant flags dropped & 12 \\
\bottomrule
\end{tabular}
\caption{Blocks of the model matrix (column counts for the default settings on the seed-0
training rows; tuning changes the vocabulary size).}
\label{tab:features}
\end{table}

Two weightings of the text block are compared. TF-IDF is the course standard. BM25~\cite{robertson2009},
the ranking function of classic search engines, differs in two ways: repeating a word has
diminishing returns (term frequency saturates, controlled by $k_1$), and long ads are discounted
against the average length (controlled by $b$). Since the two logistic regressions share every
other component, their difference isolates the effect of the word weighting.

\subsection{Metric and decision threshold}

The main metric is PR-AUC (average precision). It measures how well a model ranks frauds above
real ads without fixing a threshold, which matches the use case of a moderator queue ordered by
score. Its floor is the fraud rate, 0.049, rather than 0.5. ROC-AUC is reported as well, but it
counts the plentiful true negatives and therefore stays high even for weak
models~\cite{saito2015}: the decision tree of Section~\ref{results} has ROC-AUC 0.843
but PR-AUC 0.402.

To also report precision, recall and F1, each model's scores are turned into decisions at the
threshold that maximises F1 on the validation rows. F1 treats a missed fraud and a false alarm as
equally costly; a job board with a known cost ratio would use that ratio instead, which would move
the threshold but not the PR-AUC. The majority dummy is the only exception: it gives every ad the
same score, so a tuned threshold would flag everything, and it keeps 0.5.

\subsection{Tuning and test discipline}

Hyper-parameters are tuned by random search over the grids of Table~\ref{tab:models}: 4 candidates
per model and seed, each scored by PR-AUC with 3-fold \texttt{StratifiedGroupKFold} inside the
training rows, so the grouping also holds within cross-validation. The best candidate is refitted
on all training rows. The test rows are then scored exactly once per (model, seed), and test
numbers are not printed while experiments run, so they cannot influence any choice. Every run
saves its metrics, threshold, chosen hyper-parameters and raw scores, which is what Sections
\ref{results} and \ref{explain} are computed from.
